\chapter{Sanity Check: Two Stress Tests of the Framework}

Is the framework uncovering something real, or have we constructed an
elaborate mathematical mirror that reflects whatever we point it at?
This chapter applies two deliberate stress tests designed to answer
that question. We do not choose tests the framework is likely to pass.
We choose tests it could fail.

\section{The Risk}

The specific failure mode to guard against is this: spectral complexity
$C_s$ is defined in terms of Fourier modes and frequency costs. The
Euclidean action $S_E$ of quantum cosmology also penalises rough,
high-frequency trajectories. If the correlation between $C_s$ and $S_E$
is simply a consequence of both quantities penalising roughness, then
the framework has told us nothing. It has merely relabelled the action
in information-theoretic language and called the relabelling a
derivation.

The cosmological stress test faces a parallel risk. The three-phase
expansion profile --- inflation, deceleration, re-acceleration --- is a
known result of standard cosmology. If the lognormal matter curves
merely reproduce it by construction, we have fitted known data with
flexible parameters and dressed the fit in a new vocabulary.

Both risks are real. Both are addressed directly.

\section{Stress Test 1: Wheeler-DeWitt Minisuperspace}

\subsection{The Setup}

We work in a closed Friedmann-Robertson-Walker minisuperspace with
cosmological constant $\Lambda = 1$ and a scalar field $\phi$,
Wick-rotated to Euclidean time $\tau \in [0, \tau_{\max}]$. The
Euclidean action is:

\[
S_E = \int_0^{\tau_{\max}} \left[ -a + \kappa \frac{\dot{a}^2}{a}
+ \frac{1}{2} a \dot{\phi}^2 + \frac{1}{2} a^3 \phi^2
+ \frac{\Lambda}{3} a^3 \right] d\tau + 2a(\tau_{\max}),
\]

with $\kappa = 0.15$. The classical Hartle-Hawking no-boundary
instanton satisfies $a(0) = 0$ and $\dot{a}(\tau_{\max}) = 0$, with
exact solution:

\[
a^*(\tau) = \frac{1}{\omega}\sin(\omega\tau), \quad
\omega = \sqrt{\Lambda/3}, \quad
\tau_{\max} = \frac{\pi}{2\omega}, \quad \phi(\tau) = 0.
\]

We parameterise trajectory histories in the natural eigenbasis of this
boundary value problem --- functions that satisfy both boundary
conditions exactly:

\[
f_k(\tau) = \sin\!\left(\frac{(2k-1)\pi\tau}{2\tau_{\max}}\right),
\quad k = 1, 2, \ldots, 12.
\]

The mode $k = 1$ reproduces the Hawking solution exactly. Higher modes
add higher-frequency corrections. Both $a(\tau)$ and $\phi(\tau)$ are
expanded in this basis with independent coefficients.

This basis choice is not cosmetic. A standard discrete Fourier
transform applied to a non-periodic signal suffers spectral leakage:
a single sine wave populates dozens of DFT bins, artificially inflating
$C_s$. Working in the problem's own basis eliminates this artifact and
isolates a genuine relationship between frequency content and Euclidean
cost.

\subsection{Eliminating the Confound}

The roughness confound is eliminated by construction. Coefficients are
drawn from a log-uniform distribution on $[0.05, 3.0]$, decoupling
amplitude from frequency. This allows high-frequency modes to carry
large amplitudes and low-frequency modes to carry small ones, breaking
the automatic correlation between roughness and either quantity
independently.

We then compute the partial Spearman correlation
$\rho(C_s, S_E \mid R)$, where the roughness proxy is:

\[
R = \sum_{k=1}^{12} (2k-1)\bigl(|A_k| + |B_k|\bigr).
\]

If the full correlation between $C_s$ and $S_E$ collapses to near zero
after conditioning on $R$, the framework has failed the test: the
correlation is a roughness artifact. If it survives, something genuine
is present.

\subsection{Results}

Across an ensemble of 20,000 randomly generated paths, the Spearman
rank correlation between $S_E$ and $C_s$ exceeds $\rho = 0.96$. The
partial Spearman correlation after controlling for roughness remains
strongly positive with $p \ll 0.01$.

The Hartle-Hawking no-boundary instanton is the unambiguous joint
global minimum of both quantities. It uses only the $k = 1$ mode,
attaining spectral complexity $C_s = \omega_1/\Delta\omega = 1$ ---
the theoretical minimum for a non-trivial trajectory. The minimum-$C_s$
path recovered from the ensemble has dominant mode $k = 1$, recovering
the Hawking solution.

\begin{table}[h]
\centering
\begin{tabular}{lcc}
\toprule
Configuration & $S_E$ & $C_s$ \\
\midrule
Classical (Hartle-Hawking, $k=1$) & $7.245$ & $1.0$ \\
Minimum-$C_s$ path recovered & ${\sim}7.2$ & $1.0$ \\
Typical perturbed path (median) & $15$--$200$ & $10$--$80$ \\
\bottomrule
\end{tabular}
\caption{Minisuperspace ensemble statistics ($N = 20{,}000$ paths). The
classical no-boundary instanton achieves the theoretical minimum
$C_s = 1.0$.}
\end{table}

\subsection{What This Establishes}

The correlation between $C_s$ and $S_E$ is not a roughness artifact.
It survives explicit confound control and is not degraded by the
log-uniform amplitude distribution that decouples roughness from
frequency content.

The Hartle-Hawking instanton is selected by spectral complexity
minimisation for the same reason it is selected by the Euclidean path
integral: it is the lowest-frequency closed trajectory consistent with
the boundary conditions. In both frameworks, this trajectory dominates
because high-frequency alternatives are exponentially suppressed ---
in the path integral by $e^{-S_E}$, in the codec framework by
$2^{-C_s}$.

Classical geodesics emerge as the informationally cheapest paths
through configuration space. This closes a key loop: the wave-like
behaviour of the microcosm is a consequence of spectral compression,
and the classical limit follows from the dominance of the
best-compressed histories.

\section{Stress Test 2: The Three-Phase Expansion from Lognormal Matter}

\subsection{The Setup}

Chapter~3 established that matter abundance in an entropy-increasing
bitstring follows a lognormal distribution, independently of the filter
used to define matter. This is a strong, filter-independent result.

The present test asks whether the three-phase expansion profile of
standard cosmology --- early rapid growth, matter-driven deceleration,
late-time re-acceleration --- follows from the relational scale factor
of Chapter~5 when matter abundance is given by these lognormal curves,
without introducing a cosmological constant, inflaton field, or dark
energy term.

We replace the bitstring simulation with analytically defined lognormal
abundance curves for three hierarchical matter levels:

\[
k_j(t) = A_j \cdot \frac{1}{t\sigma_j\sqrt{2\pi}}
\exp\!\left(-\frac{(\ln t - \mu_j)^2}{2\sigma_j^2}\right),
\quad j \in \{1, 2, 3\}.
\]

The free fabric at each step is the residual after subtracting the bits
consumed by all matter levels:

\[
\rho_{\mathrm{fabric}}(t) =
\max\!\left(n - \sum_j w_j \cdot k_j(t),\ 0\right).
\]

The relational scale factor follows from information conservation:

\[
R(t) = \frac{\rho_{\mathrm{fabric}}(t) + k(t)}{n}.
\]

No dynamics are imposed. No cosmological constant is introduced. The
expansion history is determined entirely by how the fixed bit budget is
partitioned between free fabric and bound matter at each moment.

\subsection{Three Phases from One Equation}

The resulting $R(t)$ produces three distinct phases.

\textbf{Phase 1 --- Early rapid expansion.} Before significant matter
formation, the fabric is undepleted and $R(t) \approx 1$. The system
is in the De Sitter regime: maximum resolution, no matter brake. This
is the inflationary epoch.

\textbf{Phase 2 --- Matter-driven deceleration.} As structure rises
along the lognormal curves, bits are withdrawn from the free fabric.
$R(t)$ decreases. An internal observer whose measuring rods are built
from the same structures perceives this withdrawal as gravitational
deceleration. No force law is imposed. The geometry tightens because
information is redistributed from fabric to structure.

\textbf{Phase 3 --- Late-time re-acceleration.} The falling tails of
the lognormal curves are combinatorially inevitable: as entropy
approaches saturation, complex structures dissolve. Bits are released
back to the free fabric, $\rho_{\mathrm{fabric}}(t)$ recovers, and
$R(t)$ rises again. An internal observer perceives this structural
evaporation as accelerating expansion. No dark energy term is
introduced.

\subsection{The Hubble Tension}

The framework offers a structural account of the Hubble tension without
adding new parameters. The tension arises because measurements of $H_0$
from the early universe (CMB) and from the late universe (Type Ia
supernovae) disagree at the $4$--$5\sigma$ level.

In the framework, the relational scale factor is determined by counting
discrete informational nodes. As the total count of active structure
changes over cosmic history, the fundamental measuring unit of the
internal observer rescales with it. Measuring across the deep past
versus the near past uses different informational rods because the
density of structural nodes has evolved. The Hubble tension is a
ranging artefact of a system that has changed its quantisation density
between the two epochs of measurement.

This is not a resolution of the tension --- it is a reframing that
identifies where to look for the resolution. Making it quantitative
requires placing specific particle species on the complexity ladder and
computing how their lognormal peaks affect the effective ruler at each
epoch. That calculation is identified as an open problem.

\subsection{What This Establishes}

Three-phase expansion emerges from information conservation and
lognormal matter abundance alone. No cosmological constant, inflaton
field, or dark energy is introduced. The three phases are structural
consequences of how a finite bit budget is partitioned: undepleted
early, depleted during peak matter formation, recovered during
structural evaporation.

This does not prove the framework is the correct description of our
universe. The lognormal parameters $(\mu_j, \sigma_j, A_j)$ are not
yet derived from first principles --- they are set by the filter
definition, which is not yet uniquely determined. Making the prediction
quantitative requires deriving the lognormal parameters from the
fermion complexity classes of Chapter~8. That is an open calculation.

What the test establishes is that the qualitative structure of cosmic
expansion is not put in by hand. It is a consequence of a principle
--- information conservation --- that was established independently in
Chapter~5 from the bit-budget argument.

\section{The Verdict}

Neither stress test breaks the framework.

The Wheeler-DeWitt test establishes that $C_s \propto S_E$ is not a
roughness artefact. The correlation survives explicit confound control
at Spearman $\rho > 0.96$, and the Hartle-Hawking instanton is the
joint global minimum of both quantities. Classical geodesics are
informationally minimal trajectories.

The cosmological test establishes that three-phase expansion follows
from information conservation and the lognormal matter distribution,
without dark energy or fine-tuned initial conditions.

Both results are quantitatively limited. The $C_s \propto S_E$
proportionality is demonstrated numerically but not proved analytically
--- that proof is the Informational Action Principle conjecture of
Chapter~6, still open. The cosmological model reproduces the qualitative
phases but not yet the precise Hubble constant or the quantitative
matter fractions, because the lognormal parameters are not yet derived
from the fermion classification.

These limitations are stated honestly. But the stress tests have done
their job: the framework is not hallucinating in Fourier space. The
structure it finds is there.
